\documentclass[12pt,a4paper,notitlepage]{article}
\usepackage{amsmath}
\usepackage{amssymb}
\usepackage{amsbsy}
\usepackage{mathtools}
\usepackage{float}
\usepackage[french]{babel}
\usepackage{tikz}

\usepackage{palatino}

 \usepackage[active]{srcltx}
\usepackage{scrtime}

\newcounter{qnumber}
\setcounter{qnumber}{0}

\newcounter{enumber}
\setcounter{enumber}{0}

\newcommand{\question}{\textbf{(\addtocounter{qnumber}{1}\theqnumber)}\,}
\newcommand{\exercice}{\textbf{\addtocounter{enumber}{1}\textsc{Exercice} \theenumber}.\,}
\setlength{\parindent}{0cm}

% Absorbe les petits dépassements de justification (lignes trop longues)
% en répartissant l'élasticité sur les espaces plutôt que dans la marge.
\setlength{\emergencystretch}{3em}


\begin{document}

\title{\textsc{Économétrie Approfondie}\\ \small{(Session de rattrapage --- Éléments de correction)}}
\date{Juin 2026}

\maketitle

\thispagestyle{empty}


\bigskip

\exercice On suppose que les données sont générées par le modèle~:
\[
  y_i = x_{i,1}\beta_1 + \ldots + x_{i,K}\beta_K + \varepsilon_i
\]
où les $x_{i,k}$ sont des variables explicatives déterministes, les
$\beta_k$ des paramètres réels, et $\varepsilon_i$ une variable
aléatoire i.i.d. centrée de variance $\sigma_{\varepsilon}^2$.

\question On obtient la représentation matricielle en concaténant
(verticalement) les $N$ observations. On pose
$Y = \left( y_1, \dots, y_N \right)'$,
$\varepsilon = \left( \varepsilon_1, \dots, \varepsilon_N \right)'$,
$\beta = \left( \beta_1, \dots, \beta_K \right)'$ et
\[
  X =
  \begin{pmatrix}
    x_{1,1} & x_{1,2} & \ldots & x_{1,K}\\
    x_{2,1} & x_{2,2} & \ldots & x_{2,K}\\
    \vdots  & \vdots  &        & \vdots \\
    x_{N,1} & x_{N,2} & \ldots & x_{N,K}\\
  \end{pmatrix}
\]
$X$ est une matrice $N\times K$ (chaque ligne rassemble les $K$
variables explicatives d'une observation, chaque colonne les $N$
observations d'une variable). Le modèle s'écrit alors~:
\[
Y = X\beta + \varepsilon
\]

\question L'estimateur des MCO est le vecteur $\hat b$ qui minimise la
somme des carrés des résidus~:
\begin{equation*}
  \begin{split}
    \hat b &= \arg\min_{\{b\}} (Y-Xb)'(Y-Xb) \\
              &= \arg\min_{\{b\}} Y'Y - 2b'X'Y + b'X'Xb
  \end{split}
\end{equation*}
(le terme $Y'Y$ ne dépend pas de $b$ et les scalaires $b'X'Y$ et
$Y'Xb$ sont égaux). La condition du premier ordre (dérivation par
rapport au vecteur $b$) est~:
\[
2X'X\hat b - 2X'Y = 0 \Leftrightarrow \hat b = \left( X'X \right)^{-1}X'Y
\]
Pour que cet estimateur existe, il faut que la matrice $X'X$ soit
inversible, c'est-à-dire que $X$ soit de plein rang colonne
($\mathrm{rg}(X)=K$)~: absence de colinéarité parfaite entre les
variables explicatives et $N\geq K$. La condition du second ordre est
satisfaite car $X'X$ est définie positive sous l'hypothèse de plein
rang~: le point critique est bien un minimum.

\question $\hat b$ est sans biais. En substituant le DGP dans
l'expression de l'estimateur~:
\[
  \begin{split}
    \hat b &= \left( X'X \right)^{-1}X'\left( X\beta + \varepsilon \right)\\
               &= \beta + \left( X'X \right)^{-1}X'\varepsilon
  \end{split}
\]
Comme les variables explicatives sont déterministes~:
\[
  \mathbb E\left[\hat b\right] = \beta + \left( X'X \right)^{-1}X'\mathbb E\left[\varepsilon\right] = \beta
\]
puisque $\mathbb E[\varepsilon]=0$.

\question En utilisant $\hat b-\beta=\left( X'X \right)^{-1}X'\varepsilon$
et le fait que les $\varepsilon_i$ sont i.i.d.
($\mathbb V[\varepsilon]=\sigma_\varepsilon^2 I_N$)~:
\[
  \begin{split}
    \mathbb V\left[\hat b \right] &= \mathbb E\left[\left( X'X \right)^{-1}X'\varepsilon\varepsilon'X\left( X'X \right)^{-1}\right]\\
                                     &= \left( X'X \right)^{-1}X'\,\mathbb V\left[\varepsilon\right]\, X\left( X'X \right)^{-1} \\
                                     &= \sigma^2_{\varepsilon}\left( X'X \right)^{-1}X'X\left( X'X \right)^{-1}\\
                                     &= \sigma^2_{\varepsilon}\left( X'X \right)^{-1}
  \end{split}
\]
(les régresseurs étant déterministes, $\left( X'X \right)^{-1}X'$ sort
de l'espérance).

\question \textbf{Théorème de Gauss-Markov.} Sous les hypothèses du
modèle linéaire (régresseurs déterministes de plein rang, erreurs
centrées, homoscédastiques et non corrélées,
$\mathbb V[\varepsilon]=\sigma_\varepsilon^2 I_N$), l'estimateur des
MCO est, parmi tous les estimateurs linéaires et sans biais de
$\beta$, celui de variance minimale (\emph{Best Linear Unbiased
  Estimator}).

\smallskip

\textit{Démonstration.} Soit $\tilde b = CY$ un estimateur linéaire
quelconque, avec $C$ une matrice $K\times N$ déterministe. Posons
$C = (X'X)^{-1}X' + D$, où $D$ mesure l'écart aux MCO. En
substituant le DGP~:
\[
\tilde b = CY = C X\beta + C\varepsilon
\]
Pour que $\tilde b$ soit sans biais \emph{quel que soit} $\beta$, il
faut $\mathbb E[\tilde b]=CX\beta=\beta$, donc $CX=I_K$. Or~:
\[
CX = \left( (X'X)^{-1}X' + D \right)X = I_K + DX
\]
La condition de non-biais impose donc $DX = 0$. La variance de
$\tilde b$ s'écrit alors~:
\[
\mathbb V[\tilde b] = \sigma_\varepsilon^2\, CC'
\]
avec
\[
  \begin{split}
    CC' &= \left( (X'X)^{-1}X' + D \right)\left( X(X'X)^{-1} + D' \right)\\
        &= (X'X)^{-1} + (X'X)^{-1}\underbrace{X'D'}_{=(DX)'=0} \\
        &\quad + \underbrace{DX}_{=0}(X'X)^{-1} + DD'\\
        &= (X'X)^{-1} + DD'
  \end{split}
\]
d'où~:
\[
\mathbb V[\tilde b] = \sigma_\varepsilon^2 (X'X)^{-1} + \sigma_\varepsilon^2 DD'
= \mathbb V[\hat b] + \sigma_\varepsilon^2 DD'
\]
La matrice $DD'$ est semi-définie positive (pour tout vecteur $a$,
$a'DD'a = \|D'a\|^2\geq 0$). On a donc
$\mathbb V[\tilde b]-\mathbb V[\hat b]\succeq 0$~: la variance de tout
autre estimateur linéaire sans biais dépasse (au sens des matrices
semi-définies positives) celle des MCO. L'égalité n'est atteinte que
si $D=0$, c'est-à-dire $\tilde b=\hat b$. Les MCO sont donc BLUE.

\setcounter{qnumber}{0}
\bigskip

\exercice On considère le modèle générateur des données~:
\[
y_i = \beta x_i + \varepsilon_i, \qquad i=1,\ldots,N,
\]
avec $x_i$ déterministe non nul, $\mathbb E[\varepsilon_i]=0$, les
$\varepsilon_i$ indépendantes mais hétéroscédastiques~:
$\mathbb V[\varepsilon_i]=\sigma^2 w_i$, les poids $w_i>0$ étant connus.

\question En minimisant $\sum_i (y_i - b x_i)^2$, la condition du
premier ordre donne l'estimateur des MCO~:
\[
\hat b = \frac{\sum_{i=1}^N x_i y_i}{\sum_{i=1}^N x_i^2}
\]
Il est sans biais. En substituant le DGP~:
\[
\hat b = \frac{\sum_i x_i(\beta x_i + \varepsilon_i)}{\sum_i x_i^2}
       = \beta + \frac{\sum_i x_i \varepsilon_i}{\sum_i x_i^2}
\]
et comme les $x_i$ sont déterministes et $\mathbb E[\varepsilon_i]=0$~:
\[
\mathbb E[\hat b] = \beta + \frac{\sum_i x_i \mathbb E[\varepsilon_i]}{\sum_i x_i^2} = \beta
\]
L'absence de biais ne dépend pas de l'hétéroscédasticité~: elle tient
au caractère déterministe des régresseurs et à la nullité de
l'espérance des erreurs.

\question En partant de
$\hat b - \beta = \dfrac{\sum_i x_i \varepsilon_i}{\sum_i x_i^2}$ et en
utilisant l'indépendance des erreurs
($\mathbb V[\varepsilon_i]=\sigma^2 w_i$, covariances nulles)~:
\[
  \begin{split}
    \mathbb V[\hat b] &= \frac{1}{\left( \sum_i x_i^2 \right)^2}\, \mathbb V\!\left[ \sum_i x_i \varepsilon_i \right]\\
                      &= \frac{\sum_i x_i^2\, \mathbb V[\varepsilon_i]}{\left( \sum_i x_i^2 \right)^2}
                      = \frac{\sigma^2 \sum_i w_i x_i^2}{\left( \sum_i x_i^2 \right)^2}
  \end{split}
\]
On remarque que ce n'est plus la formule usuelle
$\sigma^2/\sum_i x_i^2$~: utiliser cette dernière (comme le ferait un
logiciel par défaut) conduirait à une estimation erronée de la
variance, et donc à une inférence (tests, intervalles de confiance)
invalide.

\question Pour détecter l'hétéroscédasticité, on peut utiliser un test
fondé sur les résidus des MCO $\hat\epsilon_i$~:
\begin{itemize}
\item \textbf{Test de Breusch-Pagan}~: on régresse les résidus au
  carré $\hat\epsilon_i^2$ sur les variables susceptibles d'expliquer
  la variance (ici $x_i$, $x_i^2$, ou les $w_i$). Sous l'hypothèse
  nulle d'homoscédasticité, ces variables ne doivent rien expliquer~;
  la statistique $N R^2$ de cette régression auxiliaire suit
  asymptotiquement un $\chi^2$ dont le nombre de degrés de liberté est
  le nombre de régresseurs auxiliaires. Une valeur élevée conduit à
  rejeter l'homoscédasticité.
\item \textbf{Test de White}~: même principe, en ajoutant les carrés
  et produits croisés des régresseurs (test plus général, ne
  présupposant pas la forme de l'hétéroscédasticité).
\item \textbf{Test de Goldfeld-Quandt}~: on ordonne les observations
  selon une variable soupçonnée de gouverner la variance, on estime le
  modèle séparément sur les deux sous-échantillons extrêmes, et l'on
  compare les variances résiduelles par un test de Fisher.
\end{itemize}

\question Non, l'estimateur des MCO n'est pas efficace. Le théorème de
Gauss-Markov suppose des erreurs sphériques
($\mathbb V[\varepsilon]=\sigma^2 I_N$). Ici la matrice de
variance-covariance des erreurs vaut
$\mathbb V[\varepsilon]=\sigma^2\,\mathrm{diag}(w_1,\ldots,w_N)\neq
\sigma^2 I_N$~: les MCO restent sans biais mais ne sont plus de
variance minimale. Il existe un estimateur linéaire sans biais plus
précis, l'estimateur des moindres carrés généralisés.

\question Comme $w_i>0$, on divise chaque observation par
$\sqrt{w_i}$. En posant
$\tilde y_i = y_i/\sqrt{w_i}$, $\tilde x_i = x_i/\sqrt{w_i}$ et
$\tilde\varepsilon_i = \varepsilon_i/\sqrt{w_i}$, le modèle transformé~:
\[
\tilde y_i = \beta\, \tilde x_i + \tilde\varepsilon_i
\]
a des erreurs homoscédastiques~:
\[
\mathbb V[\tilde\varepsilon_i] = \frac{1}{w_i}\,\mathbb V[\varepsilon_i] = \frac{\sigma^2 w_i}{w_i} = \sigma^2
\]
On peut alors appliquer les MCO au modèle transformé (Gauss-Markov
s'applique de nouveau)~:
\[
\hat b_{\textsc{mcg}} = \frac{\sum_i \tilde x_i \tilde y_i}{\sum_i \tilde x_i^2}
                      = \frac{\sum_i x_i y_i / w_i}{\sum_i x_i^2 / w_i}
\]
Il s'agit de l'estimateur des \emph{moindres carrés pondérés} (MCP, cas
particulier des MCG pour des erreurs hétéroscédastiques non
corrélées)~: chaque observation est pondérée par $1/w_i$, c'est-à-dire
l'inverse de sa variance. Les observations les plus «~bruitées~»
(grand $w_i$) reçoivent un poids plus faible.

\question Le modèle transformé étant homoscédastique de variance
$\sigma^2$~:
\[
\mathbb V[\hat b_{\textsc{mcg}}] = \frac{\sigma^2}{\sum_i \tilde x_i^2}
= \frac{\sigma^2}{\sum_i x_i^2 / w_i}
\]
Comparons avec $\mathbb V[\hat b]=\sigma^2 \dfrac{\sum_i w_i x_i^2}{(\sum_i x_i^2)^2}$.
En appliquant l'inégalité de Cauchy-Schwarz avec
$a_i = x_i\sqrt{w_i}$ et $c_i = x_i/\sqrt{w_i}$~:
\[
\left( \sum_i x_i^2 \right)^2 = \left( \sum_i a_i c_i \right)^2
\leq \left( \sum_i w_i x_i^2 \right)\left( \sum_i \frac{x_i^2}{w_i} \right)
\]
d'où~:
\[
  \begin{split}
    \mathbb V[\hat b] &= \sigma^2\frac{\sum_i w_i x_i^2}{\left( \sum_i x_i^2 \right)^2}
    \geq \sigma^2\frac{\sum_i w_i x_i^2}{\left( \sum_i w_i x_i^2 \right)\left( \sum_i x_i^2/w_i \right)}\\
    &= \frac{\sigma^2}{\sum_i x_i^2/w_i} = \mathbb V[\hat b_{\textsc{mcg}}]
  \end{split}
\]
On a donc bien $\mathbb V[\hat b]\geq \mathbb V[\hat b_{\textsc{mcg}}]$,
avec égalité si et seulement si les $w_i$ sont tous égaux
(homoscédasticité). Ceci confirme que les MCG/MCP sont efficaces
quand les MCO ne le sont pas.

\question Si les poids $w_i$ sont inconnus, deux stratégies~:
\begin{itemize}
\item \textbf{MCG réalisables (FGLS)}~: on spécifie une forme
  paramétrique pour la variance, $w_i = h(z_i,\theta)$ (par exemple
  $\sigma_i^2 = \sigma^2 z_i^2$). On estime d'abord le modèle par les
  MCO, on récupère les résidus $\hat\epsilon_i$, on en déduit une
  estimation $\hat\theta$ (en régressant $\ln\hat\epsilon_i^2$ sur
  $z_i$, par exemple), puis on applique les MCP avec les poids estimés
  $\hat w_i$. L'estimateur obtenu n'est plus sans biais en échantillon
  fini (les poids dépendent des données) mais reste convergent et
  asymptotiquement efficace.
\item \textbf{MCO + écarts-types robustes}~: si l'on ne souhaite pas
  modéliser la variance, on conserve l'estimateur des MCO (sans biais
  et convergent) mais on corrige l'estimation de sa variance à l'aide
  de l'estimateur robuste à l'hétéroscédasticité de White~:
  \[
  \widehat{\mathbb V}[\hat b] = \frac{\sum_i x_i^2\, \hat\epsilon_i^2}{\left( \sum_i x_i^2 \right)^2}
  \]
  ce qui rétablit la validité des tests, au prix d'une perte
  d'efficacité.
\end{itemize}

\setcounter{qnumber}{0}
\bigskip

\exercice Le (log) salaire de l'individu $i$ est généré par~:
\[
y_i = \beta_0 + \beta_1 x_i + \beta_2 z_i + \varepsilon_i
\]
où $x_i$ est le nombre d'années d'études, $z_i$ le talent (non
observé), et $\varepsilon_i$ une erreur i.i.d. centrée non corrélée à
$x_i$ et $z_i$. Les variables $x_i$ et $z_i$ sont aléatoires,
centrées, de variances $\sigma_x^2$, $\sigma_z^2$, et corrélées~:
$\sigma_{xz}=\mathrm{cov}(x_i,z_i)\neq 0$. L'économètre n'observe pas
$z_i$ et estime $y_i = b_0 + b_1 x_i + u_i$.

\question Le talent omis est rejeté dans le terme d'erreur du modèle
empirique~:
\[
u_i = \beta_2 z_i + \varepsilon_i
\]
(le modèle empirique « voit » comme une seule perturbation $u_i$ la
somme de l'effet du talent et du choc $\varepsilon_i$).

\question En utilisant la bilinéarité de la covariance, et le fait que
$x_i$ est non corrélé à $\varepsilon_i$~:
\[
  \begin{split}
    \mathrm{cov}(x_i,u_i) &= \mathrm{cov}(x_i,\beta_2 z_i + \varepsilon_i)\\
                          &= \beta_2\,\mathrm{cov}(x_i,z_i) + \mathrm{cov}(x_i,\varepsilon_i)\\
                          &= \beta_2\, \sigma_{xz} \neq 0
  \end{split}
\]
(dès lors que $\beta_2\neq 0$ et $\sigma_{xz}\neq 0$). Le régresseur
$x_i$ est donc \emph{endogène}~: il est corrélé avec l'erreur du
modèle empirique. La condition d'orthogonalité étant violée,
l'estimateur des MCO de $b_1$ est biaisé et non convergent.

\question L'estimateur des MCO de la pente s'écrit (variables
centrées)~:
\[
\hat b_1 = \frac{\frac1N\sum_i x_i y_i}{\frac1N\sum_i x_i^2}
\]
En substituant le DGP $y_i=\beta_0+\beta_1 x_i+\beta_2 z_i+\varepsilon_i$
(on travaille sur les variables centrées, la constante disparaît)~:
\[
\hat b_1 = \beta_1 + \beta_2\frac{\frac1N\sum_i x_i z_i}{\frac1N\sum_i x_i^2}
+ \frac{\frac1N\sum_i x_i \varepsilon_i}{\frac1N\sum_i x_i^2}
\]
Par la loi des grands nombres,
$\frac1N\sum_i x_i^2 \xrightarrow{p}\sigma_x^2$,
$\frac1N\sum_i x_i z_i \xrightarrow{p}\sigma_{xz}$ et
$\frac1N\sum_i x_i \varepsilon_i \xrightarrow{p}\mathrm{cov}(x_i,\varepsilon_i)=0$.
Par le théorème de Slutsky~:
\[
\text{plim}_{N\to\infty}\ \hat b_1 = \beta_1 + \beta_2\,\frac{\sigma_{xz}}{\sigma_x^2}
\]
\textbf{Interprétation.} L'estimateur des MCO ne converge pas vers
$\beta_1$~: il comporte un \emph{biais de variable omise} égal à
$\beta_2\,\sigma_{xz}/\sigma_x^2$. Le terme
$\sigma_{xz}/\sigma_x^2$ est précisément le coefficient de la
régression du talent $z$ sur les études $x$~: les MCO attribuent à
l'éducation une partie de l'effet du talent. Ici $\beta_2>0$ (le
talent augmente le salaire) et $\sigma_{xz}>0$ (les plus talentueux
étudient plus longtemps)~: le biais est \emph{positif}, on surestime
le rendement de l'éducation. Le biais disparaît si $\beta_2=0$ (le
talent n'influence pas le salaire) ou si $\sigma_{xz}=0$ (talent et
études non corrélés)~: dans ces cas l'omission est sans conséquence.

\question On dispose d'une variable $g_i$ (l'éducation des parents)
telle que $\mathrm{cov}(g_i,x_i)=\sigma_{gx}\neq 0$,
$\mathrm{cov}(g_i,z_i)=0$ et $\mathrm{cov}(g_i,\varepsilon_i)=0$, donc
$\mathrm{cov}(g_i,u_i)=\beta_2\,\mathrm{cov}(g_i,z_i)+\mathrm{cov}(g_i,\varepsilon_i)=0$.
On peut alors construire l'estimateur des \textbf{variables
  instrumentales}, en instrumentant $x_i$ par $g_i$~:
\[
\hat b_1^{\textsc{vi}} = \frac{\frac1N\sum_i (g_i-\bar g)(y_i-\bar y)}{\frac1N\sum_i (g_i-\bar g)(x_i-\bar x)}
\]
En substituant le DGP (variables centrées)~:
\[
\hat b_1^{\textsc{vi}} = \beta_1 + \frac{\frac1N\sum_i g_i u_i}{\frac1N\sum_i g_i x_i}
\]
Par la loi des grands nombres,
$\frac1N\sum_i g_i u_i \xrightarrow{p}\mathrm{cov}(g_i,u_i)=0$ et
$\frac1N\sum_i g_i x_i \xrightarrow{p}\sigma_{gx}\neq 0$, d'où par le
théorème de Slutsky~:
\[
\text{plim}_{N\to\infty}\ \hat b_1^{\textsc{vi}} = \beta_1 + \frac{0}{\sigma_{gx}} = \beta_1
\]
L'estimateur des VI est donc convergent~: l'instrument permet
d'extraire la part de la variation des études $x_i$ qui est
orthogonale au talent omis, et de restaurer l'identification du
rendement de l'éducation.

\question Un bon instrument doit satisfaire deux conditions~:
\begin{itemize}
\item \textbf{Pertinence}~: l'instrument doit être corrélé avec le
  régresseur endogène, $\mathrm{cov}(g_i,x_i)=\sigma_{gx}\neq 0$ (sinon
  le dénominateur s'annule à la limite et l'estimateur n'est pas
  défini). C'est une condition vérifiable sur les données (force de la
  première étape).
\item \textbf{Validité (exogénéité)}~: l'instrument doit être non
  corrélé avec l'erreur du modèle, $\mathrm{cov}(g_i,u_i)=0$, donc ici
  non corrélé à la fois avec le talent $z_i$ et avec $\varepsilon_i$.
  L'éducation des parents ne doit influencer le salaire de l'individu
  qu'\emph{à travers} sa propre éducation (hypothèse non testable en
  juste identification, à justifier économiquement).
\end{itemize}
Non, l'estimateur des VI n'est pas sans biais en échantillon fini~: il
s'écrit comme un rapport de variables aléatoires
$\hat b_1^{\textsc{vi}}=\beta_1 + (\frac1N\sum_i g_i u_i)/(\frac1N\sum_i g_i x_i)$,
et l'espérance d'un rapport n'est pas le rapport des espérances (le
numérateur et le dénominateur ne sont pas indépendants). Il est
seulement \emph{convergent} (sans biais asymptotiquement)~: on
renonce à l'absence de biais en petit échantillon pour corriger
l'endogénéité.

\end{document}

%%% Local Variables:
%%% mode: latex
%%% TeX-master: t
%%% End:
